\section{Introduction}
\label{sec:intro}

Renewed lunar exploration plans increase the need for readable environmental information for temperature, radiation, dust, micrometeorites, solar activity, and moonquakes \citep{nasa2020artemis}. Specialized models already exist for many of these hazards, but they usually operate as separate tools. LIPAS was built to join public data and hybrid forecasts in one decision-support interface rather than to replace flight-qualified agency systems.

This study asks three questions:
\begin{enumerate}[leftmargin=*,itemsep=0.2em]
\item[\textbf{RQ1}] Can a hybrid lunar forecasting system usefully combine machine learning and physics?
\item[\textbf{RQ2}] Does a unified multi-hazard interface improve operational usability relative to separate products?
\item[\textbf{RQ3}] Does the locked hybrid improve accuracy against selected replicated public baselines?
\end{enumerate}

Because continuous co-located surface truth is not available for every hazard, evaluation uses archived measurements, engineering-model exports, and physics or catalog proxies. That limit is carried through the results and conclusion.
